\documentclass[class=scrbook, crop=false]{standalone}
\usepackage[subpreambles=true]{standalone}
\ifstandalone
    % WARNING: Proceed with caution!

% -----------------------------------------------------------------------------------
% For package standalone
% -----------------------------------------------------------------------------------
\usepackage{import}

% -----------------------------------------------------------------------------------
% Language and typeset
% -----------------------------------------------------------------------------------
\usepackage[ngerman, english]{babel}

% Umlauts and other special characters (UTF-8)
% \usepackage[utf8]{inputenc}
\usepackage{fontspec}
\setsansfont{Arial}
% \usepackage[T1]{fontenc}  % Enable accented characters and umlauts
% LuaLatex doesn't need fontenc and uses UTF-8
% \usepackage{lmodern}  % Font face


% --------------------------------------------------------------------------------
% Page formatting
% --------------------------------------------------------------------------------
% Change the header/footer for chapter beginnings and normal pages
\usepackage[automark,headsepline]{scrlayer-scrpage}

% The package provides an easy and flexible user interface to customize the page
% layout, implementing auto-centering and auto-balancing mechanisms
% WARNING: WHEN CHANGING BCOR (Binding correction), the cover needs reworking!...
\newcommand{\theBCOR}{15mm}  % Define binding correction
\usepackage[
    bindingoffset=\theBCOR,
    % showframe, % Show boxes which indicate margins and paddings
    bottom = 3.5cm, % Margins
      left = 2.5cm,
     right = 2.5cm
] {geometry}

% The package 'float' provides a container for document objects which can not be
% broken over pages, such as tables and figures
% Needed for table and figure indexes  
\usepackage{float}

% support for landscape layout
\usepackage{lscape}

% support of \tablenotes command to add notes under table
\usepackage{threeparttable}

% To allow drawing more professional tables
\usepackage{booktabs}

% --------------------------------------------------------------------------------
% Contents
% --------------------------------------------------------------------------------
% Vector graphics (for Cover page)
\usepackage{tikz} 

% Allows additional parameters when including images
\usepackage{graphicx}

% Roman font family for all headings
\addtokomafont{disposition}{\rmfamily}

% Set the line spacing to 1.5
\usepackage[onehalfspacing]{setspace}

% Improves overall text spacing
% http://www.khirevich.com/latex/microtype/
\usepackage[stretch=10]{microtype}

% Math symbols like mu outside the math environment
\usepackage{textcomp}

% A comprehensive (SI) units package∗
% For defining SI units
\usepackage[
    range-units=single,         % Formatting ranges with single unit indication: 1 - 2 m
    range-phrase=-,             % Phrase for range: 1 - 2 m vs 1 to 2 m
    separate-uncertainty=true,  % sets +- between value and uncertainty 
    multi-part-units=repeat     % In expressions with multiple values (multi part numbers) 
                                % the unit is printed each time: 1 mm x 1 mm
] {siunitx}
% https://tex.stackexchange.com/questions/124488/multi-part-numbers-and-units-in-siunitx

% Allows Sourcecodes with highlighting 
\usepackage{listings}

% This package provides user control over the layout of the three basic list
% environments: enumerate, itemize and description
\usepackage{enumitem}
\setlist{nosep} % Remove the vertical space between \item elements in all lists

% ToDo Notes
% \setlength{\marginparwidth}{2cm}
\usepackage{todonotes}
\setuptodonotes{inline, inlinepar}
\reversemarginpar  % Put ToDo notes on the binding's side
% \usepackage{soul} % Colorful ToDo notes

% Check out colors here http://latexcolor.com/
\usepackage{xcolor}

\usepackage{amsmath}    % alignment of equations

% --------------------------------------------------------------------------------
% Other elements
% --------------------------------------------------------------------------------
% Blindtext: Organic looking text dummy
\usepackage{blindtext}

% Hyperlinks within the document (PDF)
% "hidelinks" hides visual highlighting of links
\usepackage[hidelinks]{hyperref}

% Package for Glossary and Index (Acronyms are listed in a separate list) 
\usepackage[acronym, nogroupskip]{glossaries}[=v4.49] % groupskip: alphabetic grouping of entries

\usepackage{xltabular}   % <------- FOR glossaries

% Integration and management of bibliographies
\usepackage{csquotes}   % backend=biber in biblatex needs this package
\usepackage[
    style=ieee,   % style of the bibliography, entries are sorted in alphabetic order. "ieee" is another common style.
    backend=biber,      % based on package 'biber' 
    bibencoding=utf8   % ASCII Text encoding; may use "utf8" instead
] {biblatex}

% --------------------------------------------------------------------------------
%                               PATHS & FILES
% --------------------------------------------------------------------------------
% Fix paths for standalone compiling
%\def \home {D://MS/Thesis/master-thesis}

% Package: scrlayer-scrpage
% \def \stylePath {\home/settings+/style/page}
\input{D://MS/Thesis/master-thesis/settings+/style/page}  % Load page style

% Package: graphicx
\graphicspath{{D://MS/Thesis/master-thesis/images/}}  % Set path to images

% Package: listings
\input{D://MS/Thesis/master-thesis/settings+/style/code.tex}  % Set path to style file
\lstset{inputpath={D://MS/Thesis/master-thesis/code/}} % Default path to code listings

% Package: glossaries
\input{D://MS/Thesis/master-thesis/settings+/style/symbols}  % Set path to symbols list style file
\input{D://MS/Thesis/master-thesis/settings+/style/acronyms}  % Set path to acronym list style file
\input{D://MS/Thesis/master-thesis/references/glossary_acronyms}  % Load glossary, symbol and acronyms list

% Package: biblatex
\addbibresource{D://MS/Thesis/master-thesis/references/references.bib} % Set path to bib resources

% Custom variables
\input{D://MS/Thesis/master-thesis/settings+/variables}
% --------------------------------------------------------------------------------
%                                   OPTIONAL
% --------------------------------------------------------------------------------


% Simple arithmetic for LaTeX commands
% \usepackage{calc}

% Document Elements
% -------------------

% Index
% \usepackage{imakeidx}

% compact Lists
%\usepackage{paralist}

% visual improvements for citations
% \usepackage{epigraph}

% Create pseudo code
% https://www.overleaf.com/learn/latex/Algorithms
% \usepackage{algorithm}
% \usepackage{algorithmic}
%\usepackage[noend]{algpseudocode}

% Formatting
% -------------------
% Tweaks for scrbook, redefines commands of other packages
% \usepackage{scrhack}

% Intelligent space separator (nice for superscript?)
% \usepackage{xspace}

% Allows breaks within tables
%\usepackage{tabularx}

% Allows for page breaks in tables
% \usepackage{longtable}

% allows modifying of captions
% \usepackage{caption}

% Multiline comments
%\usepackage{verbatim}

% % Custom colors
% \definecolor{dartmouthgreen}{rgb}{0.05, 0.5, 0.06}

% IF you want to define unicode characters
% \DeclareUnicodeCharacter{0229}{\c{e}}
% \DeclareUnicodeCharacter{0306}{\u{Z}}


% Document elements
% ------------------------------------

% Table package
% \usepackage{booktabs}

% Pie diagram
% \usepackage{datapie}

% Side by Side images
% \usepackage{subcaption}

% For landscape tables
%\usepackage{pdflscape}
%\usepackage{afterpage}

% Graphics can be flow around by text
%\usepackage{wrapfig}

\fi

% ----------------------------------------------------------------------------
%                               Discussion
% ----------------------------------------------------------------------------
\begin{document}
  
\chapter{Discussion}
\label{Chapter::Discussion}

This chapter presents a critical evaluation of the findings reported in Chapter~\ref{Chapter::Results}. The discussion covers key aspects such as simulation quality, empirical and data-driven creep curve representations, regression model performance, and the effectiveness of active learning. Interpretations are grounded in the context of the modeling goals and limitations inherent to the workflow.

\section{Simulation Quality and Synthetic Data}

The input parameter ranges defined for dataset generation (Table~\ref{tab:param_ranges}) were designed to broadly capture material behavior under varying stress conditions. While 3376 simulations were generated successfully during initial sampling, only 980 of 2000 additional samples converged during the active learning loop. This suggests that certain regions in the parameter space, although mathematically plausible, may lead to numerical instabilities or physically infeasible behaviors when processed by the Lagamine solver.

The discrepancy in convergence rate emphasizes the importance of convergence-aware sampling strategies, particularly in iterative learning frameworks. The use of a feasibility filter in later stages was a step in this direction but may require further refinement to enhance sampling efficiency.

\section{Empirical Model Fits}

Both empirical models, the $\theta$-projection model and the three-parameter exponential model, were evaluated for their ability to represent the simulated creep curves. The $\theta$-projection model achieved superior performance in terms of RMSE and deviation (Figures~\ref{fig:theta_fit_whole_set} and \ref{fig:theta_fit_max_deviation}) from the original creep curves, consistently producing better fits across the dataset.

The three-parameter model, while offering a simplified representation of the creep behavior, showed a slightly higher deviation (Figure~\ref{fig:three_stage_fit_whole_dataset}). This suggests that although both models are capable of capturing the general creep phases, the $\theta$-projection model is more flexible in representing nuanced variations in creep behavior.

The use of empirical models for creep prediction has been explored in earlier studies, though typically in narrower or less integrated settings. For instance, Brown and Korringa~\cite{brown2002extrapolation} established the effectiveness of the $\theta$-projection model in long-term creep extrapolation for power plant steels, showing strong agreement with experimental data. Our findings confirm this flexibility, as the $\theta$-projection model achieved the lowest RMSE across the dataset (see Table~\ref{tab:model_errors}).

In contrast, recent work by Dacol and Caetano~\cite{dacol2022modelling} applied the three-parameter exponential model for unified creep phase modeling. While computationally simple, our results indicate that this model slightly underperformed compared to the $\theta$-model in capturing creep curve shape, especially in the tertiary regimes.

\section{Dimensionality Reduction of Creep Curves}

Principal Component Analysis (PCA) was applied using three different strategies for handling non-uniform input lengths. Zero padding and normalization by rupture time led to high variance capture in fewer components, with the first 2–5 components explaining over 95\% of the total variance (Figures~\ref{fig:pca_variance_zero}, \ref{fig:pca_variance_max}, \ref{fig:pca_variance_norm}). These results suggest that the time-dependent creep behavior can be effectively compressed into low-dimensional representations, especially under standardized input transformations.

Kernel PCA (KPCA), using both RBF and polynomial kernels, showed varying reconstruction quality depending on the kernel parameters and data preprocessing. The polynomial kernel performed better under max strain padding (Figure~\ref{fig:kpca_max_fill}), but the RBF kernel performed slightly better under zero padding and normalization approach (Figures~\ref{fig:kpca_zero_fill}, \ref{fig:kpca_norm_rupture_time}). The reconstruction with normalization by rupture time captured global curve shape better.

Out of the different strategies used for handling the non-uniform input vector lengths, zero padding and the max strain padding tends to capture the artifacts and the actual creep characteristics information gets to be undermined. The normalization approach works better as it doesn't introduce any artifacts during the training stage of the models.

These findings imply that while PCA captures linear variance effectively, KPCA better handles nonlinearity at the cost of interpretability and model complexity. But these data-driven models does not seem to represent the creep curves well when it comes to the teritiary stage of creep.

\section{Supervised Regression of Material Model Parameters}

The inverse modeling task involved predicting seven Morch model parameters from simulated creep curves using four ensemble regression models. Among these, the Voting Regressor achieved the best balance of mean absolute error and $R^2$ score (Table~\ref{tab:model_performance}), indicating robustness across multiple parameter dimensions.

The parity plots (Figures~\ref{fig:RF_regression} to \ref{fig:voting_regression}) doesn't show a good agreement between predicted and true values. Some parameters, particularly those with smaller dynamic ranges, exhibited greater prediction variance, possibly due to reduced signal-to-noise ratios in the training data. The general performance of these models are quite low as is implied in the table~\ref{tab:model_performance}. The reconstructed creep curves (Figure~\ref{fig:curve_recon_model_params}) also does not seem to be well reconstructed.

Overall, ensemble models provided better generalization than individual learners, though performance gains were modest.

\section{Active Learning Efficiency}

The active learning loop aimed to iteratively explore the parameter space for physically meaningful and convergent simulations. The results (Figure~\ref{fig:pairwise_projections}) showed a progressive tightening of the feasible domain, suggesting effective exploitation of the surrogate model. However, only 980 of the 2000 queries resulted in successful simulations, indicating the model's limitations in accurately predicting convergence.

The selection strategy, based on differential evolution combined with a feasibility filter, showed promise but could benefit from integration with uncertainty quantification or probabilistic feasibility modeling to reduce the rate of failed samples. Moreover, while previous studies have used static datasets, our integration of active learning represents a novel hybrid strategy that combines physics-based solvers with intelligent sampling. Although not without convergence challenges, this approach demonstrates potential for data-efficient exploration of high-dimensional material parameter spaces.

\section{Model Performance Comparison}

The final evaluation compared all the creep curve representation methods using RMSE, minimum creep rate error, and time to rupture prediction (Table~\ref{tab:model_errors}). Both empirical models outperformed data-driven methods in terms of RMSE and creep rate error, especially the $\theta$-projection model.

However, time to rupture predictions were not possible using the empirical models, highlighting a tradeoff: empirical models provide accuracy and interpretability, while data-driven models offer generalizability and completeness.

These comparisons underline the importance of choosing representations based on the intended application—whether for interpretability, simulation compatibility, or downstream surrogate modeling.

\section{Limitations of the Approach}
\label{Section::Discussion_Limitations}

While the presented framework demonstrates promising results in creep modeling and surrogate development, several limitations remain that constrain generalizability and reliability:

\begin{itemize}
	\item \textbf{Empirical model dependence:} Both the $\theta$-projection and three-parameter models rely on empirical curve fitting. Although effective for curve compression, they lack physical interpretability with respect to underlying microstructural mechanisms, limiting their value for mechanistic insight.
	
	\item \textbf{Sensitivity to input scaling:} The supervised regression models are sensitive to input scaling and feature ranges. This may affect generalization to unseen loading regimes or materials with different baseline properties.
	
	\item \textbf{Poor inverse model accuracy:} The parity plots and performance metrics (e.g., $R^2 < 0.30$) show that inverse modeling of the full set of Morch material parameters remains inaccurate. This is likely due to parameter non-identifiability, high dimensionality, and limited signal-to-noise ratio in the input data.
	
	\item \textbf{Rupture prediction limitations:} Empirical models, while offering excellent fit, do not natively support time-to-rupture prediction, which limits their applicability for long-term life assessments.
	
	\item \textbf{Active learning convergence challenges:} Although the active learning loop improved sampling efficiency, only 980 out of 2000 queries converged. This indicates that the feasibility filter may not fully capture the solver’s stability domain, leading to wasted queries and computational cost.
	
	\item \textbf{No experimental validation:} The framework is evaluated solely on synthetic FEM-generated data. Without validation against experimental creep data, the true physical accuracy and real-world robustness of the models remain unverified.
	
\end{itemize}

\section{Sensitivity Analysis}
\label{Section::Discussion_Sensitivity}

To better understand the influence of individual input parameters on model outputs, a post hoc sensitivity analysis was performed using the trained surrogate models. This analysis was based on feature importance metrics obtained from tree-based models (e.g., Random Forest and Gradient Boosting), as well as partial dependence plots.

\subsection*{Key Findings}
\begin{itemize}
	\item \textbf{Stress (\(\sigma\)) was the most influential parameter}, consistently showing the highest feature importance across all regressors. This aligns with classical creep theory, where stress is known to exponentially affect creep rate and rupture time.
	
	\item \textbf{Hardening coefficient (\(b\)) showed high non-linear effects}, especially in early creep stages. Its sampling on a logarithmic scale helped capture sharp variations, particularly in primary creep evolution.
	
	\item \textbf{Stress exponent (\(n\)) and creep stiffness (\(K\)) were moderately influential}, primarily affecting the secondary creep rate. However, their individual contributions were often confounded by interaction effects with other parameters.
	
	\item \textbf{Softening parameters (\(S_c, s_c\)) had lower importance}, indicating that tertiary creep behavior was less reliably captured or less variable across the sampled design space.
	
	\item \textbf{Activation energy (\(Q\)) and creep coefficient (\(k_c\)) showed mixed relevance}, depending on the model target. These parameters may require reparameterization or alternative representations for better interpretability.
\end{itemize}

\subsection*{Interpretation}
The analysis reveals that the surrogate models are particularly sensitive to stress and hardening behavior, suggesting that future modeling efforts should focus on these areas for better precision. Additionally, the lower sensitivity to some softening parameters may indicate redundancy or non-identifiability under current data constraints. Incorporating global sensitivity methods such as Sobol or Morris analysis could provide further quantitative insight.




\ifstandalone
	\printglossary
	\printbibliography[heading=bibintoc]
\fi

\end{document}
